
\begin{Exercice}[10 minutes]\textbf{Maxim}
En Java, que signifie le mot-clé final appliqué à une variable?
\mcqlist{La variable ne peut pas être utilisée dans une boucle,
La variable est constante et ne peut pas changer de valeur, 
 La variable est automatiquement initialisée à zéro, 
 La variable est globale}

 \begin{solution}
Option B. 
 \end{solution}

\end{Exercice}

\begin{Exercice}[10 minutes]\textbf{Maxim}
Si on exécute le code suivant:

\begin{lstlisting}[style=pythonstyle]
x = 7 
if x % 2 == 0: 
   print("pair") 
else: 
   print("impair") 
\end{lstlisting}
Que sera affiché?

\mcqlist{impair, pair, 7, Aucun affichage}

\begin{solution}
Option A. 
\end{solution}
\end{Exercice}

\begin{Exercice}[10 minutes]\textbf{Léa}
Qu'affiche le code suivant écrit en Java?

\begin{lstlisting}[style=javastyle]
int numero_jour = 5;
switch (numero_jour) {
    case 1:                
        System.out.println("Lundi");
        break;
    case 2:
        System.out.println("Mardi");
        break;
    case 3:
        System.out.println("Mercredi");
        break;
    case 4:
        System.out.println("Jeudi");
        break;            
    case 5:                
        System.out.println("Vendredi");
    case 6:                
        System.out.println("Samedi");
        break;            
    case 7:                
        System.out.println("Dimanche");
        break;            
    default:                
        System.out.println("Ce n'est pas un jour valide !");
\end{lstlisting}

\mcqlist{Vendredi, Vendredi Samedi Dimanche, Ce n'est pas un jour valide, Vendredi Samedi}

\begin{solution}
Option D. 
\end{solution}
\end{Exercice}

\begin{Exercice}[10 minutes]\textbf{Léa}
Qu'affiche le code suivant écrit en Python?
\begin{lstlisting}[style=pythonstyle]
a = False
b = True
c = False

if (a or b) and (a or not c):
    print("output 1")
elif (a or not b) and (a or c):
    print("output 2")
if (a or not b) and (a or not c):
    print("output 3")
else:
    print("output 4")
\end{lstlisting}
\mcqlist{output 1, output 1 output 4, output 4, output 1 output 3}

\begin{solution}
Option B. 
\end{solution}
\end{Exercice}

\begin{Exercice}[10 minutes]\textbf{Aurore}
Laquelle des affirmations suivantes est correcte concernant l'expression:

\begin{equation*}
((\neg p \wedge q) \wedge (\neg p \vee \neg q)) \vee ((p \wedge q) \wedge \neg q)
\end{equation*}
\mcqlist{Si p est faux et q est vrai\, l'expression est fausse, 
Si p est vrai et q est faux\, l'expression est fausse,
Si p est faux et q est faux\, l'expression est vraie,
Toutes les affirmations sont correctes}

\begin{solution}
Option B. 
\end{solution}
\end{Exercice}

\begin{Exercice}[10 minutes]\textbf{Aurore}
On attribue à une variable \texttt{variable1} la valeur 2. On attribue à une variable \texttt{variable2} la valeur "hello". On essaie ensuite d'attribuer à \texttt{variable2} la valeur de \texttt{variable1} (en écrivant \texttt{variable2} = \texttt{variable1}). Laquelle des affirmations suivantes est correcte?

\mcqlist{
   {En Java, le code va afficher une erreur, à moins que les variables aient été les deux déclarées de type var},
   {En Python, le code va afficher une erreur},
    {En Java, le code va afficher une erreur},
   {En Python, le code va convertir l'entier deux en la chaîne de caractère "2" pour l'attribuer à variable2}
}

\begin{solution}
Option C. 
\end{solution}
\end{Exercice}

\begin{Exercice}[10 minutes]\textbf{Amina}
Qu'affiche le programme Java suivant?

\begin{lstlisting}[style=javastyle]
public class Test {
  public static void main(String[] args) {
    int a = 5;
    int b = 3;
    boolean c = (a > b) && (b > 10);
    System.out.println(c);
  }
}
\end{lstlisting}

\mcqlist{True, False, Le programme provoque une erreur de compilation, c}

\begin{solution}
Option B. 
\end{solution}

\end{Exercice}


\begin{Exercice}[10 minutes]\textbf{Amina}
Qu'affiche le programme Java suivant?
\begin{lstlisting}[style=javastyle]
public class Test {
  public static void main(String[] args) {
    int a = 5;
    int b = 3;
    boolean c = (a > b) && ((b < 10) || (a < b)) || (a == b);
    System.out.println(c);
  }
}
\end{lstlisting}
\mcqlist{False, Le programme provoque une erreur de compilation, E, True}

\begin{solution}
Option D. 
\end{solution}
\end{Exercice}

\begin{Exercice}[10 minutes]\textbf{Sean}
    Dans la table de vérité suivante, par quelle expression faut-il remplacer \#LINE REMOVED pour obtenir la troisième colonne. 
    
    \begin{center}
        \begin{tabular}{|c|c|c|}
            \hline
            $p$ & $q$ & \#LINE REMOVED \\
            \hline
            0 & 0 & 1 \\
            \hline
            0 & 1 & 1 \\
            \hline
            1 & 0 & 1 \\
            \hline
            1 & 1 & 0 \\
            \hline
        \end{tabular}
    \end{center}

    \mcqlist{
    $((\neg p \vee q) \wedge (p \vee \neg q))$,
    $(\neg ((p \wedge q) \vee (\neg p \wedge \neg q)))$,
    $((p \vee q) \wedge \neg (p \wedge q))$,
    $\neg (p \wedge q)$,
    Aucune des réponses n'est correcte
    }

\begin{solution}
Option D.
\end{solution}
\end{Exercice}
    % \item \textbf{B)} Donnez le code Python qui implémente la fonction NAND qui prend deux arguments booléens (A et B) et retourne le résultat de l'opération NAND.

    % \item \textbf{C)} Donnez une expression booléenne de la porte NAND sans utiliser l'opérateur $\text{and}$. 
\begin{Exercice}[10 minutes]\textbf{Sean}
 Soit la fonction $f(p,q) = (\neg p) \vee (\neg q)$, par quelle expression faut-il remplacer \#LINE REMOVED dans le code Java suivant pour que la méthode \texttt{computeF(boolean p, boolean q)} soit équivalente à la fonction $f$?
    \begin{lstlisting}[style=javastyle]
public static boolean computeF(boolean p, boolean q){
    return #LINE REMOVED;
}\end{lstlisting} 
\mcqlist{p \&\& (q || !q),
p || !q,
!(p \&\& q),
not p or not q,
Aucune des réponses n'est correcte
}

\begin{solution}
Option C.
\end{solution}

\end{Exercice}

\begin{Exercice}[10 minutes]\textbf{Sean}
Quelle expression est équivalente à l'expression $\neg( \neg (A \wedge A) \wedge \neg(B \wedge B))$:

\mcqlist{
     $A \lor B$, 
     $A \land B$, 
     $(\lnot A \lor \lnot B)$, 
     $\lnot(A \lor B) $, 
     Aucune des réponses n'est correcte
}

\begin{solution}
Option A (loi de Morgan). 
\end{solution}
\end{Exercice}

\begin{Exercice}[10 minutes]\textbf{Maxim}
En Python, quel opérateur réalise une division entière?

\mcqlist{
    /,
    \%,
    //,
    div()
}

\begin{solution}
Option C. 
\end{solution}
\end{Exercice}

\begin{Exercice}[5 minutes]\textbf{Maxim}
En Python, quelle exception sera levée pour une division par zéro?
\mcqlist{
    ValueError,
    TypeError,
    ZeroDivisionError,
    ArithmeticError
} 
\begin{solution}
Option C. 
\end{solution}
\end{Exercice}

\begin{Exercice}[5 minutes]\textbf{Reda}
 Soit le code python suivant qui est dans un fichier nommé \texttt{printElements.py}: 
\begin{lstlisting}[style=pythonstyle]
import sys

if __name__ == "__main__":
   arguments = sys.argv
   if len(arguments) > 1:
      for i in range(0, len(arguments) - 1):
         print(f"Elément {i+2}: {arguments[i+1]}")
   else:
      print("Assurez-vous de passer au moins un argument à votre programme") 
\end{lstlisting}


\textbf{Question 37.A} Supposons qu'on ouvre un terminal Windows dans le dossier qui contient ce fichier python, comment pouvons-nous l'exécuter tout en évitant qu'il nous imprime \texttt{Assurez-vous de passer au moins un argument à votre programme}?

\mcqlist{
    \texttt{python printElements.py one two},
    \texttt{python},
    \texttt{python HelloWorld.py},
    \texttt{python printElements.py}
}

\begin{solution}
\texttt{python printElements.py one two}
\end{solution}

\textbf{Question 37.B} Selon votre réponse à la question précédente, notez le résultat de l'exécution. 

\begin{solution}
\texttt{Elément 2: one}\\
\texttt{Elément 3: two}
\end{solution}

\textbf{Question 37.C} Quelle remarque est correcte par rapport à  la ligne 5: \texttt{if len(arguments) > 1}? 
\mcqlist{
    {Il y a une erreur, on devrait plutôt vérifier si la taille de la variable \texttt{arguments} est supérieure strictement à 0.},
    {Le nom du fichier est techniquement le premier argument, il nous faut donc plus d'un élément pour considérer que nous avons reçu des arguments.},
    {La variable \texttt{arguments} est une liste de nombres entiers.},
    {La méthode \texttt{len} est une méthode importée depuis la libraire \texttt{math}.}
}

\begin{solution}
Le nom du fichier est techniquement le premier argument, il nous faut donc plus d'un élément pour considérer que nous avons reçu des arguments.
\end{solution}

\end{Exercice}

\begin{Exercice}[5 minutes]\textbf{Reda}
Soit le code Java suivant:
\begin{lstlisting}[style=javastyle]
public class Main {
    public static void main(String[] args) {
        String secretMessage = "oatbhinepog";
        int hint = 3;
        switch(hint){
            case 1: System.out.print(secretMessage.charAt(1));
            case 10: System.out.print(secretMessage.charAt(10));
            case 6: System.out.print(secretMessage.charAt(6));
            case 3: System.out.print(secretMessage.charAt(3));
            case 50: System.out.print(secretMessage.charAt(5));
            case 60: System.out.print(secretMessage.charAt(6));
            case 100: System.out.print(secretMessage.charAt(10));
            default: System.out.print(secretMessage.charAt(0));
        }
        System.out.println();
        System.out.println(secretMessage.substring(1, 5));
        <type> value = (secretMessage.length() == 6 || secretMessage.length() == 5) && (true || hint == 3);
        System.out.println(value);
    }
}
\end{lstlisting}
Qu'est-ce qui sera affiché à la fin du bloc \texttt{switch} avant la ligne 15: \texttt{System.out.println()};

\begin{solution}
\texttt{bingo}
\end{solution}

Quelle est la valeur de \texttt{secretMessage.substring(1, 5)}?
\mcqlist{
    \texttt{"atbh"},
    \texttt{"oatb"},
    \texttt{"oatbh"},
    \texttt{"oathbhinepog"}
} 
\begin{solution}
\texttt{"atbh"}
\end{solution}


Par quoi faut-il remplacer <type>? Et quelle sera l'exécution à la ligne 17: \texttt{System.out.println(valeur);} 
\mcqlist{
    \texttt{boolean, true},
    \texttt{boolean, false},
    \texttt{String, (true || hint == 3) \&\& (secretMessage.length() == 6 || secretMessage.length() == 5) },
    \texttt{String, value}
}

\begin{solution}
\texttt{boolean, false}
\end{solution}

\end{Exercice}

\begin{Exercice}[5 minutes]\textbf{Aurore}
Parmi les éléments suivants, lesquels ne sont pas contenus dans une stack frame:
\mcqlist{
    Une adresse de retour,
    Les variables,
    Une valeur de retour,
    Les fonctions appelées,
    Tous ces éléments sont contenus dans une stack frame
}

\begin{solution}
Les fonctions appelées
\end{solution}
\end{Exercice}


\begin{Exercice}[5 minutes]\textbf{Aurore}
Qu'affiche le code suivant?
\begin{lstlisting}[style=javastyle]
public class Main {
    static void z(int a){
        System.out.print("z");
        if (a > 1){
            y(a);
        }
    }

    static void y(int a){
        System.out.print("y");
        if (a == 0){
            System.out.print(0);
        }
        else{
            a = 0;
            x(a);
        }
    }

    static void x(int a){
        System.out.print("x");
        if (a == 0) {
            y(a);
        }
        else{
            z(a);
        }
    }
    public static void main(String[] args) {
        x(2);
    }
}
\end{lstlisting}

\mcqlist{
    xz,
    xzyxy0,
    xzy0,
    Une erreur
}
\begin{solution}
xzyxy0
\end{solution}
\end{Exercice}

\begin{Exercice}[5 minutes]\textbf{Léa}
Qu'affiche le programme suivant supposant que l'utilisateur saisit l'entrée suivante:

Entrez votre username : "Algo1"
Entrez votre password : "2025-2026"

\begin{lstlisting}[style=pythonstyle]
user = "Algo"
password = "2025-2026"

try:
    user_1 = input("Entrez votre username : ")
    password_1 = input("Entrez votre mot de passe :")

    if not (user == user_1) :
        raise NameError
    else:
        if (password == password_1) :
            raise PermissionError
        else :
            print("Vous êtes connecté !")
except NameError as error :
    print("L'username n'est pas correct")
except PermissionError as error:
    print("Le mot de passe ne correspond pas à l'username")
\end{lstlisting}

\mcqlist{
 L'username n'est pas correct Le mot de passe ne correspond pas à l'username,
L'username n'est pas correct, 
Vous êtes connecté,
Le mot de passe ne correspond pas à l'username
}
\begin{solution}
L'username n'est pas correct 
\end{solution}

Maintenant qu'affiche le programme supposant que l'utilisateur saisit l'entrée suivante? \newline 

Entrez votre username : "Algo" \newline 
Entrez votre password : "2025-2027"

\mcqlist{
 L'username n'est pas correct Le mot de passe ne correspond pas à l'username,
L'username n'est pas correct, 
Vous êtes connecté,
Le mot de passe ne correspond pas à l'username
}

\begin{solution}
Vous êtes connecté
\end{solution}


\end{Exercice}

\begin{Exercice}[5 minutes]\textbf{Léa}
Qu'affiche le programme suivant?

\begin{lstlisting}[style=pythonstyle]
prix_de_groupe = 40
nb_personne = 5
prix_par_personne = prix_de_groupe / nb_personne
print(f"Le prix est de : {prix_par_personne} chf pour un groupe de {nb_personne}")
nb_personne += 1
prix_par_personne = prix_de_groupe//nb_personne
print(f"Le prix est de : {prix_par_personne} chf pour un groupe de {nb_personne}") 
\end{lstlisting}
\mcqlist{
    {Le prix est de : 8.0 chf pour un groupe de 5 \newline 
   Le prix est de : 6 chf pour un groupe de 6},
   {Le prix est de : 8 chf pour un groupe de 5 \newline 
   Le prix est de : 6 chf pour un groupe de 6},
   {Le prix est de : 8 chf pour un groupe de 5 \newline 
   Le prix est de : 6.66666666667 chf pour un groupe de 6},
   {
    Le prix est de : 8.0 chf pour un groupe de 5 \newline 
   Le prix est de : 6.66666666667 chf pour un groupe de 6
   }
}

\begin{solution}
Le prix est de : 8.0 chf pour un groupe de 5 \newline 
   Le prix est de : 6 chf pour un groupe de 6
\end{solution}
\end{Exercice}

\begin{Exercice}[5 minutes]\textbf{Léa}
Soit le programme suivant composé de 4 fonctions additions:

\begin{lstlisting}[style=javastyle]
public class Question3 {
    public static void main(String args[]) {
        System.out.println("Question3");
    }

    public static ??? addition_1 (int a, int b){
        return a + b;
    }

    public static ??? addition_2 (int a, double b){
        return a + b;
    }

    public static ??? addition_3(String a, String b){
        return a + b;
    }

    public static ??? addition_4 (double a, double b){
        System.out.println(a + b);
    }
}
\end{lstlisting}

Indiquez pour chaque fonction addition par quel type doit-on remplacer les ??? pour que le programme fonctionne (ex : addition\_5: float).

\begin{solution}
addition\_1 : int \newline 
addition\_2 : double \newline 
addition\_3 : String \newline 
addition\_4 : void
\end{solution}


\end{Exercice}

\begin{Exercice}[5 minutes]\textbf{Amina}
Qu'affiche le code suivant (Python)?
\begin{lstlisting}[style=pythonstyle]
try:
   x = int("10")
   try:
       y = int("texte")
   except ValueError:
       print("Erreur interne détectée")
   finally:
       print("Bloc interne terminé")
except Exception:
   print("Erreur externe")
finally:
    print("Bloc externe terminé")
\end{lstlisting}

\mcqlist{
    {Erreur externe \newline 
    Bloc externe terminé},
    {
        Erreur interne détectée \newline 
        Bloc interne terminé \newline 
        Bloc externe terminé
    },{
        Bloc interne terminé \newline 
        Erreur externe
    },{
        Erreur interne détectée \newline 
        Erreur externe
    }
}

\begin{solution}
Erreur interne détectée \newline 
Bloc interne terminé \newline 
Bloc externe terminé
\end{solution}

\end{Exercice}

\begin{Exercice}[5 minutes]\textbf{Amina}
Laquelle des affirmations suivantes est \textbf{incorrecte}?

\mcqlist{
    Une function en Java doit spécifier un type de retour (même void si elle ne retourne rien).,
    Le mot-clé \texttt{public} est obligatoire devant toute function.,
    Les parenthèses sont obligatoires même si la function ne prend aucun paramètre.,
    {En Java, une fonction (méthode) doit toujours être déclarée à l’intérieur d’une classe}.
}

\begin{solution}
Le mot-clé \texttt{public} est obligatoire devant toute function.
\end{solution}

\end{Exercice}


\begin{Exercice}[5 minutes]\textbf{Léa}
Qu'affiche le programme suivant?

\begin{lstlisting}[style=javastyle]
public class Main {
    public static void main(String args[]) {
        int i = 1;
        while (i % 5 != 0) {
            System.out.print(i);
            i++;
            if (i == 5) {
                i++;
            }
        }
        System.out.println(i);
    }
}
\end{lstlisting}
\mcqlist{
    1 2 3 4 5 6 7 8 9 10,
    1 2 3 4 6 7 8 9 10,
    1 2 3 4 5 6 7 8 9,
    Le programme continue en boucle
}

\begin{solution}
 1 2 3 4 6 7 8 9 10
\end{solution}

Si maintenant on rajoute une nouvelle partie au programme, le programme est maintenant comme suit: \newline 

\begin{lstlisting}[style=javastyle]
int i = 1;
while (i % 5 != 0) {
    System.out.print(i);
    i++;
    if (i == 5) {
        i ++;
    }
}
System.out.println(i);
        
// Nouvelle partie du programme 
int j = 0;
while (j % 3 == 0) {
    System.out.print(j);
    i++;
    if (j == 3) {
        j++;
    }
}
\end{lstlisting}

Qu'affiche ce programme?

\mcqlist{
    {1 2 3 4 5 6 7 8 9 10 \newline 
    0 1 2},
    {
        1 2 3 4  6 7 8 9 10 \newline 
   0
    },{
        1 2 3 4 5 6 7 8 9 \newline 
   1 2 3
    },{
        Le programme continue en boucle 
    }
}

\begin{solution}
Le programme continue en boucle
\end{solution}

\end{Exercice}


\begin{Exercice}[5 minutes]\textbf{Léa}
Quel est le résultat de la fonction suivante?

\begin{lstlisting}[style=pythonstyle]
def mystery_function(a, b):
    if b == 0:
        return 1
    else:
        return a * mystery_function(a, b - 1)
\end{lstlisting}

\mcqlist{
    a*b,
    a*(b+1),
    a**b,
    a**(b-1)
}

\begin{solution}
a**b
\end{solution}

\end{Exercice}

\begin{Exercice}[5 minutes]\textbf{Léa}
Qu'affiche le code Python suivant?

\begin{lstlisting}[style=pythonstyle]
film = {
    "titre": "Interstellar",
    "realisateur": {
        "prenom": "Christopher",
        "nom": "Nolan",
        "date" : 2014
    },
"acteurs": [
    {"prenom": "Matthew", "nom": "McConaughey", "role": "Cooper"},
    {"prenom": "Anne", "nom": "Hathaway", "role": "Brand"},
    {"prenom": "Jessica", "nom": "Chastain", "role": "Murph"}
]
}

print(film["titre"])
print(film["realisateur"]["nom"] == "Nolan")

for acteurs in film["acteurs"] :
    print(acteurs["prenom"],acteurs["nom"])
\end{lstlisting}

\mcqlist{
    Interstellar True Matthew McConaughey Anne Hathaway Jessica Chastain,
    Interstellar True Matthew Cooper Anne Brand Jessica Murph,
    Inception True Matthew McConaughey Anne Hathaway Jessica Chastain,
    Interstellar False Cooper Brand Murph
}

\begin{solution}
Interstellar True Matthew McConaughey Anne Hathaway Jessica Chastain,
\end{solution}
\end{Exercice}

\begin{Exercice}[5 minutes]\textbf{Sean}
Quelle affirmation est vraie à propos du programme Java ci-dessous:
\begin{lstlisting}[style=javastyle]
public class SyntaxeError {
    public static void main(String[] args) {
      try {
          int x = 5  // omission volontaire du ;  
      } catch (Exception e) {
          System.out.println("Erreur capturée");
          x = 5;
      }
    }
  }
\end{lstlisting}

\mcqlist{{L'erreur de syntaxe est capturée par le bloc catch au moment de l'exécution, et le programme affiche "Erreur capturée".},
{L'erreur de syntaxe ne peut pas être capturée et le programme plante au moment de l'exécution avec une SyntaxError.},
{Le programme ne compile pas à cause de l'erreur de syntaxe, donc il ne s'exécute jamais.},
{L'erreur de syntaxe est ignorée car elle est à l'intérieur d'un bloc try-catch et le programme s'exécute mais saute la déclaration de la variable.}
}
\begin{solution}
Le programme ne compile pas à cause de l'erreur de syntaxe, donc il ne s'exécute jamais.
\end{solution}
\end{Exercice}


\begin{Exercice}[5 minutes]\textbf{Sean}
Soit le code ci-dessous en Python. Par quelle expression faut-il remplacer la ligne \texttt{\#LINE REMOVED} dans la fonction \texttt{compatibilite} pour qu’elle soit équivalente à la fonction \texttt{compatibiliteRapide}?
\begin{lstlisting}[style=pythonstyle]
def compatibilite(donneur, receveur):
  # modèle simplifié par rapport à la réalité !
  donneur_groupe = donneur[0]
  donneur_rhesus = donneur[1]
  receveur_groupe = receveur[0]
  receveur_rhesus = receveur[1]

  if donneur_rhesus == receveur_rhesus:
    if donneur_groupe == "O":
      return True
    elif receveur_groupe == "AB":
      return True
    elif donneur_groupe == receveur_groupe:
      return True
  return False 

def compatibiliteRapide(donneur, receveur):
  #LINE REMOVED
\end{lstlisting}

\mcqlist{
    \texttt{return False},
    \texttt{return donneur[0] == "O" or receveur[0] == "AB" or donneur[0] == receveur[0] and donneur[1] == receveur[1]},
    \texttt{return (donneur[0] == "O" or receveur[0] == "AB" or donneur[0] == receveur[0]) or (donneur[1] == receveur[1])},
    \texttt{return (donneur[0] == "O" or receveur[0] == "AB" or donneur[0] == receveur[0]) and (donneur[1] == receveur[1])}
}

\begin{solution}
\texttt{return (donneur[0] == "O" or receveur[0] == "AB" or donneur[0] == receveur[0]) and (donneur[1] == receveur[1])}
\end{solution}

\end{Exercice}

\begin{Exercice}[5 minutes]\textbf{Sean}
Soit le code suivant en Python, qu’est-ce qui sera affiché dans la console après son exécution?

\begin{lstlisting}[style=pythonstyle]
complements = {'A': 'T', 'T': 'A', 'C': 'G', 'G': 'C'}
dna_strand = "ACTACT"
complementary_strand = ""
for base in dna_strand:
    complementary_strand += complements[base]
print(complementary_strand)
\end{lstlisting}

\mcqlist{
   ACTACT,
   TGATGA,
    AGTAGT,
 KeyError 'ACTACT'
}
\begin{solution}
TGATGA
\end{solution}
\end{Exercice}


\begin{Exercice}[5 minutes]\textbf{Sean}
Soit le code suivant en Python, qu'est-ce qui sera affiché dans la console après son exécution?

\begin{lstlisting}[style=pythonstyle]
complements = {'A': 'T', 'T': 'A', 'C': 'G', 'G': 'C'}

def comprev(dna):
    if len(dna) == 0:
        return ""
    return complements[dna[-1]] + comprev(dna[1:])

print(comprev("ACTACT"))
\end{lstlisting}

\mcqlist{
    TGATGA,
    AGTAGT,
    AAAAAA,
    RecursionError: maximum recursion depth exceeded
}

\begin{solution}
AAAAAA
\end{solution}
\end{Exercice}

\begin{Exercice}[5 minutes]\textbf{Reda}
On souhaite écrire un programme qui vérifie si une équipe peut participer à un tournoi de sport selon les critères suivants:

\begin{itemize}
\item Si l’équipe compte 3 membres, la moyenne d’âge doit être supérieure à 25.
\item Si l’équipe compte 2 membres, le ou la capitaine doit s’appeler "Alex".
\item Si l’équipe compte 1 membre, cette personne doit avoir 45 ans.
\end{itemize}

On dispose des variables suivantes:

\begin{itemize}
\item \texttt{ages}: une liste d’entiers représentant les âges des membres de l’équipe,
\item \texttt{captainName}: une chaîne de caractères contenant le nom du ou de la capitaine.
\end{itemize}

Laquelle des lignes de code Python suivantes traduit correctement ces conditions?
\begin{itemize}
  \item \texttt{(len(ages) == 3 and (ages[1] + ages[2] + ages[3])/3 > 25) or (captainName == "Alex" and len(ages) == 2) or (len(ages) == 1 and ages[1] == 45)}

  \item \texttt{(len(ages) == 3 or (ages[0] + ages[1] + ages[2])/3 > 25) and (captainName == "Alex" or len(ages) == 2) and (len(ages) == 1 or ages[0] == 45)}

  \item \texttt{(len(ages) == 3 and (ages[0] + ages[1] + ages[2])/3 > 25) or (captainName == "Alex" and len(ages) == 2) or (len(ages) == 1 and ages[0] == 45)}

  \item \texttt{(len(ages) == 3 \&\& (ages[1] + ages[2] + ages[3])/3 > 25) || (captainName == "Alex" \&\& len(ages) == 2) || (len(ages) == 1 \&\& ages[1] == 45)}
\end{itemize}

\begin{solution}
\texttt{(len(ages) == 3 and (ages[0] + ages[1] + ages[2])/3 > 25) or (captainName == "Alex" and len(ages) == 2) or (len(ages) == 1 and ages[0] == 45)}
\end{solution}


\end{Exercice}

\begin{Exercice}[5 minutes]\textbf{Reda}
Qu'affiche l'exécution du code suivant?
\begin{lstlisting}[style=pythonstyle]
def func1():
    i = 0
    value1 = 5
    value2 = 6
    def func2(value1, value2):
        global i
        i += 1
        value1 += 1
        value2 -= 1
        if i == 0:
            return value1
        else:
            return value2
    return func2(value1, value2) - func2(value1, value2)
if __name__ == "__main__":
    value1 = 3
    value2 = 4
    i = -1
    print(func1())
    i = 0
\end{lstlisting}

\mcqlist{
    5,0,1,Il y aura une erreur
}

\begin{solution}
1
\end{solution}

Supposons que l’on supprime les lignes 3, 16, 17 et 20 du code précédent.
Quelle sera la conséquence sur l’exécution du programme?

\mcqlist{
    L’exécution donnera un résultat différent de la question précédente.,
L’exécution donnera le même résultat que la question précédente.,
On ne pourra pas exécuter car le code ne compilera pas.,
Une erreur sera provoquée.,
}

\begin{solution}
L’exécution donnera le même résultat que la question précédente.
\end{solution}
\end{Exercice}

\begin{Exercice}[5 minutes]\textbf{Maxim}
En Python, qu'affiche le code suivant?
\begin{lstlisting}[style=pythonstyle]
mon_tuple = (1,2,3)
print(4 in mon_tuple)
\end{lstlisting}

\mcqlist{
    {(1, 2, 3, 4)}, True, False, {(1, 2, 3)}
}

\begin{solution}
False
\end{solution}
\end{Exercice}
\begin{Exercice}[5 minutes]\textbf{Maxim}
En Java, qu'affiche le programme suivant:

\begin{lstlisting}[style=javastyle]
import java.util.Arrays;
import java.util.List;

public class Main {
    public static void main(String[] args) {
        List<Integer> numbers = Arrays.asList(1, 2, 3, 4, 5);
        for (int x : numbers) {
           System.out.println(x);
           if (x < 3) continue;        
        }
    }
}
\end{lstlisting}
 

\mcqlist{
    1 2 3 4 5
,3 4 5
,1 2
,Rien
}
\begin{solution}
1 2 3 4 5
\end{solution}

\end{Exercice}


\begin{Exercice}[5 minutes]\textbf{Aurore}

Quelles affirmations sont correctes concernant les deux programmes suivants:

\begin{lstlisting}[style=javastyle]
List<String> bonjour = List.of("Hello","Annyeong","Hola","Hallo","Arigato","Ciao","Salaam");
        
//Programme 1
LinkedList<String> bonjour2 = new LinkedList<>(bonjour);
Iterator<String> iter = bonjour2.iterator();
while(iter.hasNext()){
    String ling = iter.next();
    if(ling.charAt(0)=='H'){
        iter.remove();
    }
}
for (int i=0;i<bonjour2.size();i++){
    System.out.println(bonjour2.get(i));
}

//Programme 2
LinkedList<String> bonjour3 = new LinkedList<>(bonjour);
for (int i=0;i<bonjour3.size();i++){
    if(bonjour.get(i).charAt(0)!='H'){
        System.out.println(bonjour3.get(i));
    }
}
\end{lstlisting}


\mcqlist{
    Les programmes afficheront la même chose,
    Le programme 1 provoque une erreur,
    Le programme 2 provoque une erreur,
    Le programme 1 modifie la liste \texttt{bonjour2},
    Le programme 2 modifie la liste \texttt{bonjour3}
}

\begin{solution}
Les programmes afficheront la même chose \newline 
Le programme 1 modifie la liste \texttt{bonjour2}

\end{solution}
\end{Exercice}

\begin{Exercice}[5 minutes]\textbf{Aurore}
Par quoi faut-il remplacer la ligne \texttt{\#LINE REMOVED} pour que le programme suivant :
\begin{lstlisting}[style=pythonstyle]
cars = {"nissan":["micra", "aryia"], "citroën":["picasso","cactus","C3"], "opel":["corsa"]}
for key in cars.keys():
    #LINE REMOVED
\end{lstlisting}
affiche la sortie suivante:

\begin{verbatim}
nissan : ['micra', 'aryia']
citroën : ['picasso', 'cactus', 'C3']
opel : ['corsa']
\end{verbatim}

\mcqlist{
    \texttt{print("{} : {}".format(key,cars[key]))},
    \texttt{print(key)},
    \texttt{print("{} : {}".format(cars[1],cars.values[key]))},
    \texttt{print("{} : {}".format("nissan", cars[key]))},
    \texttt{print(cars)}
}

\begin{solution}
\texttt{\texttt{print("{} : {}".format(key,cars[key]))}}
\end{solution}
\end{Exercice}

\begin{Exercice}[5 minutes]\textbf{Reda}

\begin{lstlisting}[style=pythonstyle]
def equivalentes(l1, l2):
      if #Q2
         return True
      if len(l1) != len(l2):
         return False
      if l1[0] not in l2:
         return False
      l2_copy = l2.copy()
      l2_copy.remove(l1[0])
      #Q3
\end{lstlisting}

Laquelle de ces propositions est fausse concernant le cas de base d’une fonction récursive?
\mcqlist{
{Le cas de base résout une partie du problème avant de poursuivre la récursion.},
{Il représente la résolution du problème dans son cas le plus simple.},
{Une fois arrivé au cas de base, nous pouvons remonter la pile (stack) d’appels pour obtenir notre réponse finale.},
Il permet dans certaines situations d’éviter un appel infini à la fonction elle-même.
}

\begin{solution}
Le cas de base résout une partie du problème avant de poursuivre la récursion.
\end{solution}

Par quoi devons-nous remplacer \#Q2 pour bien représenter le cas de base dans notre contexte?

\mcqlist{
    \texttt{liste1[0] == liste2[0]:},
    \texttt{len(liste1) == 0 and len(liste2) == 0:},
    \texttt{liste1 == liste2:},
    \texttt{len(liste1) == 0 or len(liste2) == 0:}
}
\begin{solution}
    \texttt{len(liste1) == 0 and len(liste2) == 0:}
\end{solution}

Par quoi devons-nous remplacer \#Q3 pour que notre programme soit correct?
\mcqlist{
    \texttt{return equivalentes(liste1, liste2\_copy)},
    \texttt{return equivalentes(liste1[1:], liste2\_copy)},
    \texttt{return equivalentes(liste1[1:], liste2)},
    \texttt{return equivalentes(liste1[:len(l1)], liste2[:len(l2)])}
}
\begin{solution}
    \texttt{return equivalentes(liste1[1:], liste2\_copy)}
\end{solution}

Imaginons que nous souhaitons vérifier si deux listes de nombres entiers sont équivalentes par itération. Laquelle de ces stratégies ne fonctionnerait pas ?
(On suppose que les deux listes ont la même taille et que chacune des stratégies peut être correctement implémentée.)

\mcqlist{
    {Pour chaque élément de liste1, on vérifie si cet élément est contenu dans liste2.},
    {On crée un dictionnaire qui associe à chaque élément de liste1 son nombre d’apparitions ; pour chaque élément de liste2, on vérifie que c’est une clé du dictionnaire et que la valeur correspond à son nombre d’apparitions dans liste2.},
    {Pour chaque élément de liste1, on vérifie s’il apparaît le même nombre de fois dans liste1 et liste2.},
    {On réarrange chaque liste dans l’ordre croissant, puis on compare indice par indice du début à la fin.}
}
\begin{solution}
Pour chaque élément de liste1, on vérifie si cet élément est contenu dans liste2.
\end{solution}

\end{Exercice}

\begin{Exercice}[5 minutes]\textbf{Amina}

Qu'affiche le code Python suivant?

\begin{lstlisting}[style=pythonstyle]
numbers = [1, 2, 3, 4, 5, 6, 7, 8]
for x in numbers:
    if x < 6:
        numbers.remove(x)
print(numbers)
\end{lstlisting}

\mcqlist{
    {$[$6, 7, 8$]$},
    {$[$2, 4, 6, 7, 8$]$},
    {$[$1, 3, 5, 6, 7, 8$]$},
    RuntimeError: Concurrent modification error
}

\begin{solution}
[2, 4, 6, 7, 8]
\end{solution}


\end{Exercice}

\begin{Exercice}[5 minutes]\textbf{Amina}

On souhaite supprimer toutes les entrées dont la valeur est inférieure à 0. La variable suivante est donnée :

\begin{lstlisting}[style=javastyle]
Map<String, Integer> scores = new HashMap<>();
scores.put("a", 3);
scores.put("b", -1);
scores.put("c", 0);
scores.put("d", -5);
\end{lstlisting}

Lesquelles de ces versions ne provoque pas d’exception et produit le résultat attendu ?


\begin{lstlisting}[style=javastyle]
// Option A
for (String key : scores.keySet()) {
    if (scores.get(key) < 0) {
        scores.remove(key);
    }
}
System.out.println(scores);
\end{lstlisting}

\begin{lstlisting}[style=javastyle]
// Option B
for (Map.Entry<String, Integer> e : scores.entrySet()) {
    if (e.getValue() < 0) {
        scores.remove(e.getKey());
    }
}
System.out.println(scores);
\end{lstlisting}

\begin{lstlisting}[style=javastyle]
// Option C
Iterator<Map.Entry<String, Integer>> it = scores.entrySet().iterator();
while (it.hasNext()) {
    if (it.next().getValue() < 0) {
        it.remove();
    }
}
System.out.println(scores);
\end{lstlisting}

\begin{lstlisting}[style=javastyle]
// Option D
Iterator<String> it = scores.keySet().iterator();
while (it.hasNext()) {
    String k = it.next();
    if (scores.get(k) < 0) {
        scores.remove(k);
    }
}
System.out.println(scores);
\end{lstlisting}

\begin{solution}
C
\end{solution}

\end{Exercice}

\begin{Exercice}[5 minutes]\textbf{Amina}
Qu'affiche le code Python suivant?

\begin{lstlisting}[style=pythonstyle]
def mystery(n, acc=[]):
    if n == 0:
        return acc
    acc.append(n)
    return mystery(n - 1)

a = mystery(3)
b = mystery(2)
print(a, b)
\end{lstlisting}

\mcqlist{
{$[$3, 2, 1$]$ $[$2, 1$]$},
{$[$3, 2, 1$]$ $[$3, 2, 1$]$},
{$[$3, 2, 1, 2, 1$]$ $[$3, 2, 1, 2, 1$]$},
Provoque une erreur: \texttt{RecursionError}
}

\begin{solution}
Option C.
\end{solution}

\end{Exercice}

